\chapter{Reconfiguración v2: requisitos, contenedor y liberación del drogue}\label{cap:arq}
\usepgfplotslibrary{groupplots}

El anexo A cambió, y el equipo aclaró dos libertades de diseño: el drogue puede soltarse
antes de abrir el segundo sistema de descenso, y el anexo no exige una pared si el
CanSat va dentro de un contenedor que elige el equipo. Este capítulo reúne lo que eso cambia
en el modelo aerodinámico y en el diseño. El resultado principal es que, con el drogue
soltado, la etapa 2 es el rotor solo, y el radio y la masa que vuela con el rotor pasan a
ser las dos palancas de diseño: $V_2\propto\sqrt{m_2}/R$. Con $R=\SI{0.25}{m}$ el rotor
solo da $V_2=\SI{6.4}{m/s}$ ($\langle V\rangle=\SI{7.1}{m/s}$ sobre todo el descenso). Si la
organización permite separar el contenedor con el drogue, la masa que vuela baja a unos
\SI{0.36}{kg} y $V_2$ a \SI{5.5}{m/s}. Los demás capítulos (BEMT, arranque, estabilidad,
estructura) se recalcularon con estos datos.

\begin{resultado}[Resumen del capítulo]
\begin{itemize}
  \item \textbf{El promedio de \Req{C4} importa.} Leído sobre todo el descenso desde el
        apogeo, $\langle V\rangle=1/(0{,}2/V_1+0{,}8/V_2)$ y exige $V_2\le\SI{7.26}{m/s}$. El
        rotor solo de $R=\SI{0.20}{m}$ ($V_2=\SI{7.85}{m/s}$) no lo cumple; el de
        $R=\SI{0.25}{m}$ sí ($\SI{7.11}{m/s}$), con un margen de \SI{0.9}{m/s}.
  \item \textbf{Dos arquitecturas con el mismo rotor.} La Arquitectura~1 mantiene el cuerpo
        impreso (Ø\,81) con las palas exteriores y suelta el drogue con un pasador, después de
        confirmar que el anillo de traba giró. La Arquitectura~2 usa un contenedor tubular no
        impreso que se separa con el drogue: la carga útil cae por el extremo abierto en
        \SI{0.13}{s}, las palas quedan libres y la masa del contenedor no vuela con el
        rotor ($m_2\approx\SI{0.36}{kg}$).
  \item \textbf{Con la Arquitectura~2:} $V_2=\SI{5.5}{m/s}$, $\langle V\rangle=\SI{6.2}{m/s}$ y
        \SI{5}{J} de energía de aterrizaje (contra \SI{10}{J}); el perfil de la pala puede
        tener espesor ($t/c$ hasta 11\,\% con chasis Ø\,72). Requiere que la organización
        permita separar el contenedor: es una consulta, no un hecho.
  \item \textbf{Integración de la Arquitectura~1} (sección~\ref{sec:arq-integracion}): cuatro aletas
        antigiro de $40\times\SI{25}{mm}$ en la banda inferior de \SI{25}{mm} (calaje \ang{1.7},
        \SI{84}{\percent} del amortiguamiento de las de $40\times40$); copa del drogue de
        Ø\,75 $\times$ \SI{10}{mm} llena al \SI{59}{\percent} (estimado; probar con maqueta); la masa
        sube a \SI{497}{g} y deja \SI{3}{g} de lastre: pesar el panel solar decide si hay riesgo de
        sobrepeso (\SI{44}{\percent} con \SI{37}{g}, \SI{1.9}{\percent} con \SI{8}{g}), y la segunda celda no es viable.
  \item \textbf{Simulación de cumplimiento} (Monte Carlo, $2\times10^5$ corridas, con el apogeo
        de diseño): con la Arquitectura~1 y $R=\SI{0.25}{m}$, \Req{C4} se cumple como régimen en el
        \SI{99.7}{\percent} de los casos y como promedio global en el \SI{86}{\percent};
        con la Arquitectura~2, en el \SI{100}{\percent} y el \SI{99.8}{\percent}.
  \item \textbf{Recomendación.} Fijar $R=\SI{0.25}{m}$ (rotor común) y llevar adelante la
        Arquitectura~1 como base, con la Arquitectura~2 lista para activarse si la organización
        responde que sí. La matriz de decisión da ventaja a la 2 con el \SI{81}{\percent} de
        probabilidad frente a variaciones de los pesos; el factor decisivo es el riesgo
        normativo.
\end{itemize}
\end{resultado}

%=====================================================================
\section{Qué cambió y qué hay que consultar}\label{sec:arq-cambios}

\begin{table}[H]
\centering\small
\caption{Cambios del anexo~A (v2) y de las condiciones de diseño, y su efecto.}\label{tab:arq-cambios}
\begin{tabularx}{\textwidth}{@{}lYYY@{}}
\toprule
\textbf{Tema} & \textbf{Antes (v1)} & \textbf{Ahora (v2)} & \textbf{Efecto}\\\midrule
\Req{S2} diámetro & \SI{95}{mm} (las bases dicen $95\times200$) & \SI{90}{mm} & Plegado rehecho; cuerpo Ø\,81 (sec.~\ref{sec:arq-geom}); inconsistencia a consultar\\
\Req{E7} autonomía & --- & mínimo 2\,h integrado en el cohete & Cámaras al despegue, no en la rampa (sec.~\ref{sec:dim-prom} y cap.~\ref{cap:sens})\\
\Req{X3} & --- & comando automático al separarse del cohete & El firmware detecta la separación (cap.~\ref{cap:electronica})\\
\Req{X5} & aceleraciones & aceleraciones y rotaciones en 3 ejes, GPS & MPU6050 cubre el giroscopio (cap.~\ref{cap:sensores})\\
Trama y estados & definidos en el anexo & a criterio del equipo & Se mantienen los estados propios del equipo\\
\Req{S6} pared & $\ge\SI{2}{mm}$ & $\ge\SI{2}{mm}$ \emph{si se imprime} & No hay pared obligatoria en un contenedor no impreso\\
Drogue & atado, cuerda por el eje & puede soltarse antes de abrir el rotor & Etapa 2 con el rotor solo\\
\bottomrule
\end{tabularx}
\end{table}

Las aclaraciones del equipo tienen tres consecuencias de diseño. Primera, la etapa 2 deja de
tener el aporte del drogue y su estela, y pierde el \SI{2}{\percent} de masa del drogue y
su cuerda; el rotor solo con $R=\SI{0.20}{m}$ no alcanza (sección~\ref{sec:dim-prom}). Segunda, si
el contenedor es nuestro, se puede elegir un tubo delgado no impreso, y las palas ya no
necesitan arrollarse sobre un cuerpo de pared de \SI{2}{mm}: caben con espesor de
perfil. Tercera, si además el contenedor puede separarse con el drogue, la masa que lleva
el rotor se reduce.

\paragraph{Consultas a la organización.}
\begin{enumerate}[leftmargin=*]
  \item \Req{S2}: ¿el diámetro es \SI{90}{mm} (anexo v2) o $95\times200$ (bases y condiciones)?
  \item \Req{C4}: ¿el promedio de $5\pm\SI{3}{m/s}$ se mide sobre todo el descenso desde el
        apogeo o solo desde la liberación del segundo sistema?
  \item ¿Se puede soltar el drogue antes de desplegar el segundo sistema, de modo que caiga
        aparte del CanSat?
  \item ¿Se puede separar el contenedor del CanSat en vuelo (con el drogue)? La masa de
        \Req{S1} ¿se verifica combinada solo en el pesaje?
  \item \Req{S6}: ¿la pared mínima de \SI{2}{mm} aplica solo a piezas impresas en 3D, y se
        admite un contenedor tubular no impreso (aluminio, fibra de vidrio, PVC)?
  \item Altura sobre el mar del sitio de lanzamiento y apogeo esperado. El equipo estima la
        liberación del rotor a \SI{500}{m}, un apogeo de \SI{625}{m} con el \SI{80}{\percent} de
        \Req{C4} (sección~\ref{sec:dim-tiempos}), por confirmar. La densidad del aire es la
        segunda variable de $V_2$ (capítulo~\ref{cap:sens}).
  \item \Req{X3}: ¿«separación» es la salida del CanSat del cohete?
  \item \Req{C4}: ¿la liberación es el \SI{80}{\percent} del apogeo medido o una altura fija sobre
        el suelo? Un umbral fijo reintroduce el error de escala de la atmósfera estándar
        (capítulo~\ref{cap:sens}).
\end{enumerate}

%=====================================================================
\section{Dos arquitecturas con el mismo rotor}\label{sec:arq-dos}

El rotor es común: $R=\SI{0.25}{m}$, pala de \SI{206}{mm}, $B=4$, $c=\SI{45}{mm}$,
$\theta=\ang{-3}$, bisagra con tope elástico a \ang{+15}. Las arquitecturas difieren en el
cuerpo, la liberación y la masa que vuela.

\begin{figure}[H]
\centering
\begin{tikzpicture}[x=0.40mm, y=0.40mm]
  % ================= Arquitectura 1
  \begin{scope}
    \node[font=\small\bfseries] at (0,270) {Arquitectura 1};
    % cuerpo
    \filldraw[fijo] (-40.5,25) rectangle (40.5,225);
    % espuma
    \filldraw[pattern=north east lines, pattern color=cgris] (-40.5,0) rectangle (40.5,25);
    \draw[fijo] (-40.5,0) rectangle (40.5,25);
    \node[font=\scriptsize, cgris, fill=white, inner sep=1pt] at (0,12) {espuma 25};
    % palas plegadas
    \filldraw[rot] (-44.5,25) rectangle (-40.5,231);
    \filldraw[rot] (40.5,25) rectangle (44.5,231);
    % anillo de traba
    \draw[act, line width=1.2pt] (-38,40) -- (38,40);
    \node[font=\scriptsize, cfreno, anchor=west] at (46,40) {anillo de traba (SG90)};
    % bateria y electronica
    \filldraw[fill=white, draw=cgris] (-9,60) rectangle (9,125);
    \node[font=\scriptsize, cgris] at (0,92) {18650};
    \filldraw[fill=white, draw=cgris] (-32,135) rectangle (32,205);
    \node[font=\scriptsize, cgris, align=center] at (0,170) {electrónica\\y cámaras};
    % cubo
    \filldraw[rot] (-22,225) rectangle (22,240);
    \draw[eje] (0,215) -- (0,250);
    % copa del drogue y pasador
    \filldraw[cont, fill=ccont!15] (-37.5,240) rectangle (37.5,250);
    \draw[cuerda, line width=1.4pt] (0,243) -- (0,262) node[above, font=\scriptsize, ccuerda] {drogue};
    \filldraw[act] (-8,244) rectangle (8,247);
    \node[font=\scriptsize, cfreno, anchor=west] at (46,246) {pasador (SG90)};
    \draw[cota] (-45,-12) -- (45,-12) node[midway, fill=white, font=\scriptsize] {Ø\,90 (S2)};
    \draw[cota] (-58,0) -- (-58,250) node[midway, fill=white, font=\scriptsize, rotate=90] {250 (S3)};
  \end{scope}
  % ================= Arquitectura 2
  \begin{scope}[shift={(230,0)}]
    \node[font=\small\bfseries] at (0,270) {Arquitectura 2};
    % contenedor: tubo
    \draw[cont, line width=1.0pt] (-45,0) -- (-45,240);
    \draw[cont, line width=1.0pt] (45,0) -- (45,240);
    \filldraw[cont, fill=ccont!25] (-45,240) rectangle (45,250);
    \node[font=\scriptsize, ccont, anchor=west] at (50,120) {contenedor (tubo)};
    % chasis
    \draw[fijo] (-36,0) rectangle (36,226);
    \filldraw[pattern=north east lines, pattern color=cgris] (-36,0) rectangle (36,22);
    % palas plegadas en el espacio anular
    \filldraw[rot] (-44,22) rectangle (-37,226);
    \filldraw[rot] (37,22) rectangle (44,226);
    % pestillo
    \filldraw[act] (-6,226) rectangle (6,240);
    \node[font=\scriptsize, cfreno, anchor=west] at (50,233) {pestillo (SG90)};
    % cubo
    \filldraw[rot] (-18,212) rectangle (18,226);
    % electrónica
    \filldraw[fill=white, draw=cgris] (-9,55) rectangle (9,120);
    \node[font=\scriptsize, cgris] at (0,88) {18650};
    \filldraw[fill=white, draw=cgris] (-30,130) rectangle (30,200);
    \node[font=\scriptsize, cgris, align=center] at (0,165) {electrónica\\y cámaras};
    \draw[cuerda, line width=1.4pt] (0,250) -- (0,262) node[above, font=\scriptsize, ccuerda] {drogue};
    \draw[flecha, cfreno] (0,-8) -- (0,-35) node[below, font=\scriptsize, cfreno, align=center] {la carga útil\\sale por abajo};
    \draw[cota] (-45,-52) -- (45,-52) node[midway, fill=white, font=\scriptsize] {Ø\,90 (S2)};
  \end{scope}
\end{tikzpicture}
\caption{Las dos arquitecturas, plegadas (cortes longitudinales, medidas en mm). En la
primera, las palas (naranja) cuelgan por fuera del cuerpo impreso y un anillo interno las
retiene; el drogue sale de la copa al soltar un pasador. En la segunda, las palas se pliegan
en el espacio anular entre el chasis y un tubo no impreso; al soltar el pestillo, el
drogue arrastra el contenedor hacia arriba y la carga útil cae por el extremo abierto.}
\label{fig:arq-dos}
\end{figure}

\begin{table}[H]
\centering\small
\caption{Comparación de las dos arquitecturas.}\label{tab:arq-comparacion}
\begin{tabularx}{\textwidth}{@{}lYY@{}}
\toprule
 & \textbf{Arquitectura 1} & \textbf{Arquitectura 2}\\\midrule
Cuerpo & PETG Ø\,81, pared \SI{2}{mm} (\Req{S6}) & Chasis esqueleto + tubo no impreso Ø\,90\\
Palas & Placa de \SI{0.5}{mm} arrollada por fuera & Perfil con espesor, en el espacio anular\\
Traba de las palas & Anillo interno con servo + pasador del drogue (2 servos) & El propio contenedor; un pestillo (el segundo servo es redundante)\\
Drogue & Se suelta con un pasador tras confirmar el anillo & Atado al contenedor, que se separa\\
Masa en la etapa 2 & \SI{0.49}{kg} & $\approx\SI{0.36}{kg}$\\
$V_2$, $\langle V\rangle$ ($R=\SI{0.25}{m}$) & \SI{6.4}{m/s}, \SI{7.1}{m/s} & \SI{5.5}{m/s}, \SI{6.2}{m/s}\\
Energía de aterrizaje & \SI{10}{J} & \SI{5}{J}\\
Dependencia normativa & Ninguna adicional & Separar el contenedor en vuelo\\
Largo & En el límite de \Req{S3} (copa aplanada) & Tubo de \SIrange{235}{245}{mm}\\
Antenas (GPS, LoRa) & Cuerpo plástico: transparente & Contenedor de fibra de vidrio o PVC: transparente\\
\bottomrule
\end{tabularx}
\end{table}

%=====================================================================
\section{Geometría: plegado en Ø\,90 y espacio para un perfil con espesor}\label{sec:arq-geom}

\paragraph{Arquitectura 1.} Una pala de cuerda $c$ y flecha $s=0{,}07c$, apoyada por el centro
sobre un cuerpo de radio $\rho$ con la cara cóncava hacia él, tiene el borde exterior a
\begin{equation}
  r_\text{borde}=\sqrt{(\rho-s+t)^2+(c/2)^2},
  \label{eq:arq-rborde}
\end{equation}
con $t$ el espesor. Para $\rho=\SI{40.5}{mm}$, $c=\SI{45}{mm}$ y $t=\SI{0.5}{mm}$,
$r_\text{borde}=\SI{44.03}{mm}$: el envolvente es Ø\,88,1, con \SI{1.0}{mm} de juego
radial hasta Ø\,90. El cuerpo no puede pasar de $\rho=\SI{41.6}{mm}$ (Ø\,83). Con cuerda
de \SI{50}{mm} el límite baja a \SI{40.4}{mm} (Ø\,80,8) y con \SI{40}{mm} sube a
\SI{42.6}{mm}. La figura~\ref{fig:arq-geom}(a) muestra la relación.

\paragraph{Arquitectura 2.} Con un tubo de Ø\,90 y pared de \SI{0.8}{mm}, el radio interior
es \SI{44.2}{mm}. Si el chasis tiene radio $\rho_c$, el espesor máximo del perfil que cabe
plegado es
\begin{equation}
  t_\text{máx}=\sqrt{R_i^2-(c/2)^2}-(\rho_c-s).
  \label{eq:arq-tmax}
\end{equation}
Con un chasis de Ø\,72 ($\rho_c=\SI{36}{mm}$) y $c=\SI{45}{mm}$, $t_\text{máx}=\SI{5.2}{mm}$ (11,5\,\% de la
cuerda); con Ø\,76, \SI{3.2}{mm} (7\,\%); con Ø\,68, \SI{7.2}{mm}
(figura~\ref{fig:arq-geom}(b)). Un perfil de 3 a 6\,\% de espesor entra con holgura. El
espesor sube la rigidez a torsión de la pala como $t^3$, que es lo que falta a la pala
de \SI{206}{mm} (capítulo~\ref{cap:estructura}), y permite un perfil de carbono sobre núcleo de
espuma o un perfil impreso hueco. El efecto aerodinámico se estudió en la
sección~\ref{sec:polares-espesor-v2}: una placa gruesa de espesor constante suma hasta \SI{0.5}{m/s} a la caída, pero un perfil con espesor
NACA de \SI{8}{\percent} mantiene $V_2=\SI{6.3}{m/s}$ con rpm un 8\,\% menores.

\begin{figure}[H]
\centering
\begin{tikzpicture}
\begin{groupplot}[group style={group size=2 by 1, horizontal sep=1.8cm},
  width=0.47\textwidth, height=5.6cm, grid=major, legend style={font=\scriptsize}]
\nextgroupplot[xlabel={Radio del cuerpo $\rho$ [mm]}, ylabel={Radio del borde $r_\text{borde}$ [mm]},
  title={\footnotesize (a) Arquitectura 1}, xmin=36, xmax=44, ymin=40, ymax=48, legend pos=north west]
\addplot[cfijo, thick] table[x=rho, y=c40] {analysis/data/v2_geom_env.dat};
\addlegendentry{$c=\SI{40}{mm}$}
\addplot[crot, very thick] table[x=rho, y=c45] {analysis/data/v2_geom_env.dat};
\addlegendentry{$c=\SI{45}{mm}$}
\addplot[ccuerda, thick] table[x=rho, y=c50] {analysis/data/v2_geom_env.dat};
\addlegendentry{$c=\SI{50}{mm}$}
\addplot[cfreno, dashed, thick, domain=36:44] {45};
\node[font=\scriptsize, cfreno, anchor=south east] at (axis cs:43.8,45.1) {Ø\,90};
\addplot[only marks, mark=*, mark size=2.5pt, cfreno] coordinates {(40.5,44.03)};
\nextgroupplot[xlabel={Radio del chasis $\rho_c$ [mm]}, ylabel={$t_\text{máx}$ [mm]},
  title={\footnotesize (b) Arquitectura 2}, xmin=30, xmax=41, ymin=0, ymax=12, legend pos=north east]
\addplot[cfijo, thick] table[x=rho, y=t_c40_w08] {analysis/data/v2_geom_anillo.dat};
\addlegendentry{$c=\SI{40}{mm}$}
\addplot[crot, very thick] table[x=rho, y=t_c45_w08] {analysis/data/v2_geom_anillo.dat};
\addlegendentry{$c=\SI{45}{mm}$}
\addplot[ccuerda, thick] table[x=rho, y=t_c50_w08] {analysis/data/v2_geom_anillo.dat};
\addlegendentry{$c=\SI{50}{mm}$}
\addplot[cgris, dashed] coordinates {(30,2.7) (41,2.7)};
\node[font=\scriptsize, cgris, anchor=south west] at (axis cs:30.2,2.8) {6\,\% de $c=45$};
\end{groupplot}
\end{tikzpicture}
\caption{(a) Radio del borde de la pala plegada frente al radio del cuerpo (ec.~\ref{eq:arq-rborde},
$t=\SI{0.5}{mm}$). (b) Espesor máximo del perfil en el espacio anular del contenedor
(ec.~\ref{eq:arq-tmax}, pared de \SI{0.8}{mm}).}
\label{fig:arq-geom}
\end{figure}

\paragraph{Largo.} Con $R=\SI{0.25}{m}$ las palas cuelgan \SI{206}{mm} desde las bisagras; en la
Arquitectura~1 eso obliga a poner las bisagras a \SI{231}{mm}, la copa del drogue a
\SI{10}{mm} y el total en \SI{250}{mm}, el límite de \Req{S3} (tabla~\ref{tab:largo}). Es
el punto más ajustado de esa arquitectura. En la Arquitectura~2 el contenedor mide
\SIrange{235}{245}{mm}.

%=====================================================================
\section{Contenedor, masa y energía}\label{sec:arq-masa}

En la Arquitectura~2 el contenedor es un tubo no impreso. La elección del material
decide la masa que se va con él, la transparencia a las antenas y el costo.

\begin{table}[H]
\centering\small
\caption{Tubos candidatos para el contenedor (Ø\,90 $\times$ \SI{225}{mm}; masa del tubo solo,
sin tapa; densidades típicas).}\label{tab:arq-tubos}
\begin{tabularx}{\textwidth}{@{}lccYY@{}}
\toprule
\textbf{Material} & \textbf{Pared} [mm] & \textbf{Masa} [g] & \textbf{RF (GPS, LoRa)} & \textbf{Observaciones}\\\midrule
Aluminio 6061 & 0,8 & 137 & Bloquea: antenas fuera o ventana & Mecanizable; liviano para su rigidez\\
Fibra de vidrio (GRP) & 0,6 & 73 & Transparente & Liviano; el lastre va en la tapa\\
Carbono & 0,5 & 49 & Bloquea & Rígido; costo alto\\
PVC (caño) & 1,8 & 160 & Transparente & Barato y pesado: útil como lastre que se va\\
Cartón técnico & 1,5 & 76 & Transparente & Débil al aplastamiento; sensible a la humedad\\
\bottomrule
\end{tabularx}
\end{table}

Como el contenedor se va con el drogue, un contenedor más pesado \emph{baja} la masa que
vuela con el rotor, siempre que el total siga en $500\pm\SI{10}{g}$ (\Req{S1}). Eso
favorece un tubo transparente a las antenas (GRP o PVC) más un lastre de acero o plomo en la
tapa: la telemetría debe funcionar durante la etapa 1, con el contenedor cerrado
(\Req{X3}). Un tubo metálico o de carbono exigiría una ventana o una antena exterior.

La tabla~\ref{tab:masa} del capítulo~\ref{cap:dim} da los presupuestos de ambas
arquitecturas con la electrónica real del equipo (\SI{177}{g}) y los cambios de la
sección~\ref{sec:arq-integracion}: la carga útil de la Arquitectura~2 pesa \SI{352}{g} y el
resto, hasta \SI{500}{g}, lo pone el contenedor (de fibra de vidrio, \SI{93}{g}), el
drogue y el lastre que se van. Con $m_2=\SI{0.36}{kg}$ (el valor del modelo, que cubre los
\SI{352}{g} con su desviación) y $R=\SI{0.25}{m}$ resulta
$V_2=\SI{5.5}{m/s}$ y una energía de aterrizaje de \SI{5.4}{J}.

\paragraph{Sensibilidad de la masa.} El panel solar (\SI{37}{g} en la planilla, dato dudoso)
aporta el \SI{65}{\percent} de la varianza de la masa de la carga útil en la
Arquitectura~2 (\SI{56}{\percent} en la 1); le siguen el cableado (14 y 12\,\%) y, en la 1,
el cuerpo impreso (16\,\%). En la Arquitectura~1, la masa en la etapa 2 queda fijada
por el pesaje ($500\pm10$, menos el drogue) porque el lastre se ajusta al final; en la
Arquitectura~2 depende de la suma de los componentes de la carga útil, con una desviación de
\SI{15}{g} que cambia $V_2$ en solo \SI{0.12}{m/s} (capítulo~\ref{cap:sens}).

%=====================================================================
%=====================================================================
\section{Integración de la Arquitectura 1: aletas, copa y masa}\label{sec:arq-integracion}

Tres cosas que los capítulos anteriores trataron por separado tienen que caber a la vez en el
mismo cuerpo de Ø\,81 $\times$ \SI{250}{mm}: las cuatro aletas antigiro (capítulo~\ref{cap:estabilidad}),
la copa del drogue (sección~\ref{sec:arq-secuencia}) y los cambios de electrónica del
capítulo~\ref{cap:sensores}. Los cálculos están en \texttt{analysis/v2\_integracion.py}
(salida \texttt{analysis/data/v2\_integracion.json}). Lo que depende de la geometría exacta
se marca «a confirmar en CAD»: no hay un modelo CAD todavía.

\subsection{Aletas en la banda inferior}
Las palas plegadas cuelgan por fuera del cuerpo desde las bisagras (\SI{231}{mm}) hasta las
pestañas (\SI{25}{mm}); entre $z=0$ y \SI{25}{mm} el cuerpo no tiene palas. Esa banda ya está
ocupada por el anillo de espuma de aterrizaje (Ø\,81/58 $\times$ \SI{25}{mm}, capítulo~\ref{cap:estructura}),
y es el único lugar para las aletas. Cada una es una placa de \SI{0.5}{mm} (carbono o FR-4, \SI{0.78}{g})
con una bisagra de eje vertical en la raíz y un resorte de torsión; plegada, se enrolla tangente al cuerpo en un
bolsillo de \SI{1.5}{mm} de profundidad del anillo, dentro de Ø\,90 (radio $40{,}5+0{,}5+1{,}5<\SI{45}{mm}$), y su arco
ocupa \SI{40}{mm} de los \SI{254}{mm} de perímetro (cuatro aletas: \SI{160}{mm}). El anillo de traba
suma cuatro dientes de alambre más, que bajan hasta la banda y retienen el borde libre de cada
aleta; cuando el anillo gira \ang{15} para liberar las palas, los dientes recorren \SI{10}{mm}
de arco, más que los \SI{3}{mm} de apoyo, y las aletas se abren por resorte hasta un tope
regulable, que fija el calaje de \ang{1.7}. Sin sensor propio: la confirmación del anillo (fin de carrera o Hall) sirve
a las cuatro. La carga sobre la bisagra al abrirse es mínima (presión dinámica de \SI{110}{Pa} a
\SI{13.4}{m/s} sobre \SI{1000}{mm^2}: \SI{0.1}{N}).

\begin{table}[H]
\centering\small
\caption{Aletas antigiro de la Arquitectura~1: cuatro unidades, centro a $r_f=45+s/2$ (mm).
$c$ es el amortiguamiento; $\omega_\text{eq}$, el giro de equilibrio sin calaje; $\delta$, el
calaje que anula el par de los rodamientos ($M_f=\SI{0.65}{mN.m}$).}\label{tab:arq-aletas}
\begin{tabular}{@{}lcccccc@{}}
\toprule
\textbf{Envergadura $\times$ cuerda} [mm] & $\mathcal A$ & $a_f$ & $S$ [mm$^2$] & $c$ [N.m.s] & $\omega_\text{eq}$ [rad/s] & $\delta$ [\si{\degree}]\\\midrule
$40\times40$ (no cabe en la banda)  & 2,0 & 2,60 & 1600 & $2{,}74\times10^{-4}$ & 2,4 & 1,4\\
\textbf{$40\times25$ (adoptada)}    & 3,2 & 3,48 & 1000 & $2{,}29\times10^{-4}$ & 2,8 & 1,7\\
$50\times25$ (sobresale de Ø\,90 al plegarla) & 4,0 & 3,88 & 1250 & $3{,}71\times10^{-4}$ & 1,8 & 1,1\\
\bottomrule
\end{tabular}
\end{table}

Con la cuerda limitada a \SI{25}{mm}, la aleta pierde el \SI{16}{\percent} del amortiguamiento
respecto de la de $40\times40$ y el calaje sube de \ang{1.4} a \ang{1.7}; el efecto sobre el
criterio de video ($\le\SI{0.5}{rad/s}$) es el mismo que ya tenía el diseño: depende de ajustar el calaje
en un ensayo de caída. Las cuatro aletas con resortes y los dientes pesan \SI{7}{g}
(\SI{1.6}{g} por aleta y \SI{0.4}{g} los dientes).

\subsection{Copa del drogue}
La copa de Ø\,75 $\times$ \SI{10}{mm} sobre el cubo tiene \SI{44.2}{cm^3} y, con paredes de \SI{1.5}{mm} y
fondo de \SI{1}{mm}, \SI{36.6}{cm^3} útiles. Dentro van el drogue de ripstop de \SI{40}{g/m^2}
(Ø\,\SI{250}{mm}, \SI{2.6}{g} con costuras), ocho cordones de \SI{0.25}{m} y el destorcedor, y el carro con el
pasador de liberación. Empaquetando la tela a \SI{0.35}{g/cm^3} (valor típico de un pliegue
manual, a medir) el volumen ocupado es

\begin{center}
\begin{tabular}{@{}lr@{}}
\toprule
Tela y cordones empaquetados & \SI{10.7}{cm^3}\\
Destorcedor & \SI{3.0}{cm^3}\\
Carro y pasador & \SI{8.0}{cm^3}\\\midrule
\textbf{Total} y llenado de los \SI{36.6}{cm^3} & \textbf{\SI{21.7}{cm^3}, 59\,\%}\\
\bottomrule
\end{tabular}
\end{center}

El llenado de $59$\,\% es viable pero justo: son estimaciones sin maqueta. La prueba es empaquetar el
drogue real en una copa impresa antes de congelar el CAD (lista de ensayos del
capítulo~\ref{cap:ensayos}); si el llenado pasa del \SI{80}{\percent}, la salida es comprimir
\SI{15}{mm} la pila interna (de \SI{90}{mm} a \SI{75}{mm} de PCB y cámaras, a confirmar en CAD) y alargar la copa a
\SI{25}{mm}, con lo que el llenado baja a \SI{22}{\percent}. Bajar las bisagras 15 mm (a \SI{216}{mm}) para ganar
altura no sirve: las puntas de las palas bajarían a $z=\SI{10}{mm}$, dentro de la banda de las aletas.

\subsection{Presupuesto de masa}
Los tres cambios que se suman a la tabla~\ref{tab:masa} son: \SI{+16}{g} de electrónica
(capítulo~\ref{cap:sensores}), \SI{+7}{g} por corregir la masa de la cabeza del rotor
(\SI{71}{g} contra \SI{64}{g}) y \SI{+7}{g} de las aletas. El subtotal de la Arquitectura~1 sube de \SI{467}{g} a
\SI{497}{g} y el lastre baja de \SI{33}{g} a \SI{3}{g}. Con las desviaciones de la tabla y un
pesaje final uniforme entre \SI{490}{g} y \SI{510}{g}, un Monte Carlo (\num{200000} casos)
da lo siguiente.

\begin{table}[H]
\centering\small
\caption{Presupuesto de masa de la Arquitectura~1 (Monte Carlo): subtotal sin lastre,
lastre necesario y probabilidades.}\label{tab:arq-masa-mc}
\begin{tabular}{@{}lcccc@{}}
\toprule
\textbf{Caso} & \textbf{Subtotal} [g] & \textbf{Lastre} [g] & $P(\text{lastre}<0)$ & $P(\text{lastre}<\SI{10}{g})$\\\midrule
Tabla anterior (sin cambios)                     & $467\pm20$ & 33  & 6\,\%  & 13\,\%\\
Con cambios, panel de la planilla (\SI{37}{g})   & $497\pm21$ & 3   & \textbf{44}\,\% & 63\,\%\\
Con cambios, panel de \SI{8}{g} (a pesar)        & $468\pm14$ & 32  & 1,9\,\% & 7,7\,\%\\
Con cambios y segunda celda, panel de \SI{37}{g} & $545\pm21$ & $-45$ & 98\,\% & 99,5\,\%\\
Con cambios y segunda celda, panel de \SI{8}{g}  & $516\pm14$ & $-16$ & 85\,\% & 95\,\%\\
\bottomrule
\end{tabular}
\end{table}

La conclusión es doble. Primero, \textbf{pesar el panel solar decide el presupuesto}: con los
\SI{37}{g} de la planilla la Arquitectura~1 queda sin lastre y hay que sacar masa en casi la mitad de los casos; con un panel de \SI{8}{g}
vuelve a tener el margen de antes. Segundo, \textbf{la segunda celda 18650 (\SI{48}{g}) no es viable} en la
Arquitectura~1 aunque el panel sea liviano: con el panel de \SI{8}{g} sigue con un \SI{85}{\percent} de
sobrepeso. Si el panel pesa mucho, las palancas son el cuerpo impreso (\SI{130}{g}, el
\SI{26}{\percent} del total) y el cableado; no se tocó la pared mínima de \SI{2}{mm} de \Req{S6}. En la
Arquitectura~2 la masa de la carga útil con los mismos cambios es de \SI{352}{g}; el contenedor
tiene que ser de fibra de vidrio (\SI{93}{g} con su tapa) y no de aluminio (\SI{157}{g}), y el lastre de
\SI{45}{g} viaja dentro y se suelta con él.

\section{Mapa aerodinámico: radio, masa y promedio de \texorpdfstring{\Req{C4}}{C4}}\label{sec:arq-mapa}

Con la BEMT en equilibrio de fuerzas ($T+D_\text{cuerpo}=W_2$) y par nulo, para
$B=4$, $c=\SI{45}{mm}$, $\theta=\ang{-3}$:

\begin{figure}[H]
\centering
\begin{tikzpicture}
\begin{groupplot}[group style={group size=2 by 1, horizontal sep=1.8cm},
  width=0.47\textwidth, height=6.4cm, grid=major, xlabel={Masa en la etapa 2 $m_2$ [kg]},
  xmin=0.30, xmax=0.51, legend style={font=\scriptsize}]
\nextgroupplot[ylabel={$V_2$ de régimen [m/s]}, ymin=4.8, ymax=8.4, legend pos=north west]
\addplot[cfijo, thick, mark=square*, mark size=1.2pt] table[x=m2, y=V_R20] {analysis/data/v2_mapa_ancho.dat};
\addlegendentry{$R=\SI{0.20}{m}$}
\addplot[ccuerda, thick, mark=triangle*, mark size=1.5pt] table[x=m2, y=V_R23] {analysis/data/v2_mapa_ancho.dat};
\addlegendentry{$R=\SI{0.23}{m}$}
\addplot[crot, very thick, mark=*, mark size=1.2pt] table[x=m2, y=V_R25] {analysis/data/v2_mapa_ancho.dat};
\addlegendentry{$R=\SI{0.25}{m}$}
\addplot[cfreno, dotted, thick, domain=0.30:0.51] {8};
\addplot[cgris, dotted, thick, domain=0.30:0.51] {7.26};
\node[font=\scriptsize, cgris, anchor=north west] at (axis cs:0.31,7.22) {7,26};
\nextgroupplot[ylabel={Promedio global $\langle V\rangle$ [m/s]}, ymin=5.4, ymax=9.2, legend pos=north west]
\addplot[cfijo, thick, mark=square*, mark size=1.2pt] table[x=m2, y=Vg_R20] {analysis/data/v2_mapa_ancho.dat};
\addplot[ccuerda, thick, mark=triangle*, mark size=1.5pt] table[x=m2, y=Vg_R23] {analysis/data/v2_mapa_ancho.dat};
\addplot[crot, very thick, mark=*, mark size=1.2pt] table[x=m2, y=Vg_R25] {analysis/data/v2_mapa_ancho.dat};
\addplot[cfreno, dotted, thick, domain=0.30:0.51] {8};
\node[font=\scriptsize, cfreno, anchor=south west] at (axis cs:0.31,8.05) {límite 8};
\draw[cgris, dashed] (axis cs:0.49,5.4) -- (axis cs:0.49,9.2);
\draw[cgris, dashed] (axis cs:0.36,5.4) -- (axis cs:0.36,9.2);
\node[font=\scriptsize, cgris, rotate=90, anchor=south] at (axis cs:0.49,5.6) {Arq.~1};
\node[font=\scriptsize, cgris, rotate=90, anchor=south] at (axis cs:0.36,5.6) {Arq.~2};
\end{groupplot}
\end{tikzpicture}
\caption{Velocidad de régimen y promedio global de \Req{C4} frente a la masa que vuela en la
etapa 2, para tres radios (BEMT, rotor solo, $V_1=\SI{13.44}{m/s}$). Las líneas
punteadas son el límite de \SI{8}{m/s} y, en la izquierda, el $V_2$ máximo que da
$\langle V\rangle=\SI{8}{m/s}$.}
\label{fig:arq-mapa}
\end{figure}

\begin{table}[H]
\centering\small
\caption{Rotor solo (BEMT) para tres radios y las dos masas de las arquitecturas.}\label{tab:arq-mapa}
\begin{tabular}{@{}lcccccc@{}}
\toprule
$R$ [m] & $m_2$ [kg] & $V_2$ [m/s] & rpm & $\langle V\rangle$ [m/s] & $E_\text{aterr.}$ [J] & Cumple lectura global\\\midrule
0,20 & 0,49 & 7,85 & \num{1692} & 8,56 & 15,1 & No\\
0,23 & 0,49 & 6,88 & \num{1387} & 7,63 & 11,6 & Sí (margen \SI{0.4}{m/s})\\
\textbf{0,25} & \textbf{0,49} & \textbf{6,36} & \textbf{\num{1229}} & \textbf{7,11} & \textbf{9,9} & \textbf{Sí (\SI{0.9}{m/s})}\\
0,20 & 0,33 & 6,47 & \num{1408} & 7,22 & 6,9 & Sí\\
0,23 & 0,33 & 5,67 & \num{1150} & 6,41 & 5,3 & Sí\\
0,25 & 0,33 & 5,24 & \num{1017} & 5,97 & 4,5 & Sí\\
\bottomrule
\end{tabular}
\end{table}

\paragraph{Cómo se llega al valor de diseño.}
\begin{itemize}
  \item $R=\SI{0.25}{m}$ es el menor radio con el que la Arquitectura~1 cumple ambas lecturas
        de \Req{C4} con margen. Con $R=\SI{0.23}{m}$ el promedio global queda a
        \SI{0.4}{m/s} del límite, menos que la incertidumbre de la BEMT.
  \item El radio de \SI{0.25}{m} cabe: la pala de \SI{206}{mm} ocupa el largo del cuerpo y
        el largo total queda en el límite de \Req{S3}.
  \item Con la Arquitectura~2, $R=\SI{0.23}{m}$ ya da $\langle V\rangle=\SI{6.4}{m/s}$ y reduciría
        la pala a \SI{186}{mm}. Se mantiene \SI{0.25}{m} para no cambiar el rotor entre
        arquitecturas.
  \item Subir el diámetro del drogue o cambiar $V_1$ casi no mueve el promedio global: de
        $V_1=\SI{12}{m/s}$ a \SI{18}{m/s} el promedio varía \SI{0.1}{m/s}
        (figura~\ref{fig:arq-prom}).
\end{itemize}

\begin{figure}[H]
\centering
\begin{tikzpicture}
\begin{axis}[width=0.8\textwidth, height=6cm, grid=major, xlabel={$V_2$ de régimen [m/s]},
  ylabel={Promedio global $\langle V\rangle$ [m/s]}, xmin=4, xmax=9, ymin=4, ymax=11,
  legend style={font=\scriptsize, at={(0.03,0.97)}, anchor=north west}]
\addplot[cfijo, thick] table[x=V2, y=Vg_V12] {analysis/data/v2_promedio.dat};
\addlegendentry{$V_1=\SI{12}{m/s}$}
\addplot[crot, very thick] table[x=V2, y=Vg_V134] {analysis/data/v2_promedio.dat};
\addlegendentry{$V_1=\SI{13.44}{m/s}$}
\addplot[ccuerda, thick] table[x=V2, y=Vg_V15] {analysis/data/v2_promedio.dat};
\addlegendentry{$V_1=\SI{15}{m/s}$}
\addplot[cfreno, thick, dashed] table[x=V2, y=Vg_V18] {analysis/data/v2_promedio.dat};
\addlegendentry{$V_1=\SI{18}{m/s}$}
\addplot[cfreno, dotted, thick, domain=4:9] {8};
\end{axis}
\end{tikzpicture}
\caption{Promedio global de \Req{C4} (ec.~\ref{eq:promedio}) en función de $V_2$ para varias
velocidades de la etapa 1.}
\label{fig:arq-prom}
\end{figure}

\paragraph{Tiempos de descenso.} Con el apogeo de diseño (sección~\ref{sec:dim-tiempos}), la
etapa~1 dura \SI{9.3}{s} y la etapa~2, \SI{78.6}{s} con la Arquitectura~1
($V_2=\SI{6.36}{m/s}$) o \SI{91.4}{s} con la 2 ($V_2=\SI{5.47}{m/s}$): \SI{87.9}{s} y
\SI{100.7}{s} de descenso. Si el apogeo fuera de \SI{500}{m} (liberación a \SI{400}{m}),
serían \SI{70.3}{s} y \SI{80.6}{s}.

%=====================================================================
\section{Secuencia de liberación y soltado del drogue}\label{sec:arq-secuencia}

\subsection{Caída libre al soltar}
Al soltar el drogue el cuerpo acelera hasta que las palas frenan. Con el cuerpo de Ø\,90
($C_D=0{,}8$) la velocidad terminal sin rotor es \SI{39}{m/s} ($m=\SI{0.49}{kg}$). La
tabla~\ref{tab:arq-libre} da lo que se gana en cada retardo.

\begin{table}[H]
\centering\small
\caption{Caída libre desde $V_1=\SI{13.44}{m/s}$ sin drogue ni rotor.}\label{tab:arq-libre}
\begin{tabular}{@{}lccccc@{}}
\toprule
Retardo [s] & 0,1 & 0,2 & 0,3 & 0,5 & 1,0\\\midrule
Velocidad ($m=\SI{0.49}{kg}$) [m/s] & 14,3 & 15,1 & 16,0 & 17,6 & 21,3\\
Caída [m] & 1,4 & 2,9 & 4,4 & 7,8 & 17,5\\
Velocidad ($m=\SI{0.36}{kg}$) [m/s] & 14,3 & 15,0 & 15,8 & 17,3 & 20,6\\
\bottomrule
\end{tabular}
\end{table}

Soltar el drogue \emph{antes} de que las palas empiecen a abrir es peligroso: si las palas
no abren (traba o anillo atascado), el CanSat acelera hacia \SI{39}{m/s} y no hay segundo
sistema que lo frene. El rotor, en cambio, frena sin problema desde \SIrange{14}{16}{m/s}
(capítulo~\ref{cap:arranque}). Por eso la secuencia lleva un enclavamiento: el drogue
\emph{no} se suelta hasta que el mecanismo de las palas confirme su liberación.

\subsection{Arquitectura 1: anillo, confirmación y pasador}
\begin{enumerate}[leftmargin=*]
  \item $t=0$: el servo del anillo gira \ang{15} (el SG90 gira \ang{60} en \SI{0.1}{s}, por lo
        que son $\approx\SI{0.03}{s}$). Las palas quedan libres y empiezan a abrirse por el
        resorte y el aire.
  \item $t\approx\SI{0.1}{s}$: un fin de carrera o el sensor hall confirma que el anillo giró.
        Si no confirma en \SI{0.3}{s}, se reintenta dos veces y, si falla, el drogue
        \textbf{no} se suelta (modo degradado seguro: el CanSat sigue a \SI{13}{m/s}, cumple
        \Req{C3} pero no \Req{C4}).
  \item $t\approx\SI{0.15}{s}$: el segundo servo retira el pasador y el lazo de la cuerda se
        libera. La tensión del drogue a $V_1$ es \SI{4.3}{N}; el drogue se va hacia arriba
        con una aceleración relativa muy alta y cae a $\approx\SI{2}{m/s}$.
  \item $t\approx\SI{0.3}{s}$: las palas terminan de abrir ($\SI{0.18}{s}$ de apertura) a
        \SIrange{14}{15}{m/s}; el rotor arranca (cap.~\ref{cap:arranque}).
\end{enumerate}
El pasador retirado con un servo sufre una fuerza de extracción $\mu F$ ($\mu=0{,}2$--0,4)
con $F=\SI{4.3}{N}$ en régimen: un par de \SIrange{10}{21}{mN.m} en un horn de \SI{12}{mm},
contra los \SI{180}{mN.m} del SG90~\cite{arq-sg90} (margen 8,6 a 17). En la apertura del
drogue (hasta \SI{23}{N} a \SI{25}{m/s}) el par subiría a \SIrange{55}{110}{mN.m}, con margen
1,6 a 3,3: por eso el pasador va en un carro con retén que toma la carga de apertura y el
servo solo lo mueve cuando la tensión es la de régimen (\Req{M3}). El anillo de traba usa el
segundo servo. El capítulo~\ref{cap:sensores} recomienda el MG90S de engranajes metálicos (misma huella, \SI{13}{g}) para los dos, porque el
engranaje de plástico del SG90 es el eslabón débil frente al choque (\Req{M3}); el par nominal es el mismo. Ningún mecanismo usa calor ni pirotecnia (\Req{M1}, \Req{M2}).

\subsection{Arquitectura 2: separación del contenedor}
Un pestillo (hoz o bola retenida) sujeta la carga útil al contenedor. Al soltarlo, el drogue
frena el contenedor ($D_d=\SI{4.3}{N}$ contra su peso de $\SI{1.45}{N}$) y la carga útil, más
pesada, cae por el extremo abierto. La integración de las dos aceleraciones da
(figura~\ref{fig:arq-sabot}):
\begin{itemize}
  \item la carga útil sale del tubo ($\SI{230}{mm}$ de recorrido relativo) en \SI{0.13}{s}, con el
        contenedor a \SI{11.3}{m/s} y la carga útil a \SI{14.5}{m/s};
  \item un contenedor más liviano sale antes (\SI{0.11}{s} con \SI{0.10}{kg}) y uno más
        pesado después (\SI{0.18}{s} con \SI{0.25}{kg});
  \item las palas empiezan a abrirse cuando salen del tubo; con \SI{0.18}{s} de apertura el
        rotor está abierto a los $\approx\SI{0.31}{s}$ del pestillo, a \SIrange{14}{15}{m/s}.
\end{itemize}

\begin{figure}[H]
\centering
\begin{tikzpicture}
\begin{groupplot}[group style={group size=2 by 1, horizontal sep=1.8cm},
  width=0.47\textwidth, height=5.4cm, grid=major, xlabel={Tiempo desde el pestillo [s]}, legend style={font=\scriptsize}]
\nextgroupplot[ylabel={Velocidad [m/s]}, legend pos=north west, xmin=0, xmax=0.15]
\addplot[cfijo, very thick] table[x=t, y=vp] {analysis/data/v2_sabot.dat};
\addlegendentry{carga útil}
\addplot[ccont, very thick] table[x=t, y=vc] {analysis/data/v2_sabot.dat};
\addlegendentry{contenedor + drogue}
\nextgroupplot[ylabel={Recorrido relativo [mm]}, xmin=0, xmax=0.15, legend pos=north west]
\addplot[crot, very thick] table[x=t, y expr=\thisrow{z}*1000] {analysis/data/v2_sabot.dat};
\addplot[cfreno, dashed, domain=0:0.15] {230};
\node[font=\scriptsize, cfreno, anchor=south east] at (axis cs:0.148,235) {largo del tubo};
\end{groupplot}
\end{tikzpicture}
\caption{Separación del contenedor (Arquitectura~2): velocidades y recorrido relativo de la
carga útil de \SI{0.35}{kg} dentro del contenedor de \SI{0.15}{kg} (con drogue y lastre).}
\label{fig:arq-sabot}
\end{figure}

\subsection{Cronología}
\begin{figure}[H]
\centering
\begin{tikzpicture}[x=118mm, y=1cm]
  \draw[->] (0,0) -- (0.62,0) node[right, font=\scriptsize] {$t$ [s]};
  \foreach \x in {0,0.1,0.2,0.3,0.4,0.5}{\draw (\x,0.07) -- (\x,-0.07) node[below, font=\scriptsize] {\x};}
  % arquitectura 1
  \node[font=\scriptsize\bfseries, anchor=east] at (-0.01,2.2) {Arq. 1};
  \fill[cfreno!50] (0,2.0) rectangle (0.03,2.4);          \node[font=\tiny, anchor=west] at (0.03,2.2) {anillo};
  \fill[cfijo!40] (0.03,1.5) rectangle (0.10,1.9);         \node[font=\tiny, anchor=west] at (0.10,1.7) {confirmación};
  \fill[cfreno!50] (0.13,1.0) rectangle (0.20,1.4);        \node[font=\tiny, anchor=west] at (0.20,1.2) {pasador: se va el drogue};
  \fill[crot!50] (0.03,0.5) rectangle (0.21,0.9);          \node[font=\tiny, anchor=west] at (0.21,0.7) {palas abren (0,18 s)};
  % arquitectura 2
  \node[font=\scriptsize\bfseries, anchor=east] at (-0.01,-1.0) {Arq. 2};
  \fill[cfreno!50] (0,-1.2) rectangle (0.08,-0.8);         \node[font=\tiny, anchor=west] at (0.08,-1.0) {pestillo};
  \fill[ccont!40] (0.08,-1.7) rectangle (0.23,-1.3);       \node[font=\tiny, anchor=west] at (0.23,-1.5) {sale del tubo (0,13 s)};
  \fill[crot!50] (0.23,-2.2) rectangle (0.41,-1.8);        \node[font=\tiny, anchor=west] at (0.41,-2.0) {palas abren};
\end{tikzpicture}
\caption{Cronología de la liberación, de \SI{0}{} a \SI{0.5}{s}. Desde que las palas están
abiertas, el arranque del rotor dura \SI{1.5}{s} más y el CanSat cae \SI{21.7}{m}, el
\SI{4.3}{\percent} de la altura de la etapa~2 (capítulo~\ref{cap:arranque}).}
\label{fig:arq-cronologia}
\end{figure}

\subsection{Modos de falla de la liberación}
\begin{table}[H]
\centering\small
\caption{Fallas de la liberación y su efecto.}\label{tab:arq-fallas}
\begin{tabularx}{\textwidth}{@{}lYYY@{}}
\toprule
\textbf{Falla} & \textbf{Arq.~1} & \textbf{Arq.~2} & \textbf{Consecuencia}\\\midrule
No libera la traba & Anillo trabado: no se suelta el drogue (enclavamiento) & Pestillo trabado: sigue todo junto a \SI{13}{m/s} & Cumple \Req{C3}, no \Req{C4}; seguro\\
No suelta el drogue & Pasador trabado: rotor con drogue atado, $V_2=\SI{5.8}{m/s}$ & (se separa con el contenedor) & Cumple \Req{C4}; riesgo de enredo con la cuerda\\
Palas abren y no arrancan & Rotor parado: cae a \SI{14}{m/s} & Ídem & No cumple \Req{C4}; sin freno extra\\
Contenedor roza las palas & --- & Tubo o rieles frenan las palas al salir & Apertura lenta; ensayo de caída\\
Liberación prematura (lógica) & Mismo árbol que en v1 (cap.~\ref{cap:confiabilidad}) & Ídem & Votación 2 de 3 y respaldo por tiempo\\
\bottomrule
\end{tabularx}
\end{table}

%=====================================================================
\section{Cumplimiento de C3 y C4: simulación Monte Carlo}\label{sec:arq-mc}

Se simularon $2\times10^5$ descensos por caso con variación de la masa de pesaje
(490--510~g), la masa en la etapa 2, la densidad del aire (temperatura de 5 a 35\,°C y altura del
sitio de 0 a \SI{800}{m}), $C_D$ del drogue (0,75--0,90), el tiempo de transición de $V_1$ a
$V_2$ (2,5 a \SI{5}{s}) y una incertidumbre relativa del \SI{12}{\percent} en $C_R$ (la de la
BEMT). El apogeo es fijo, el de diseño (sección~\ref{sec:dim-tiempos}). La transición se modela
como $t_s$ segundos a la velocidad media $(V_1+V_2)/2$; el rango es más ancho que los
\SI{2.6}{s} en que la simulación del capítulo~\ref{cap:arranque} deja la velocidad a menos del
\SI{2}{\percent} de $V_2$ (el $t_{90}=\SI{1.55}{s}$ es el giro del rotor, no la velocidad). La
densidad se evalúa a la altitud del sitio y no a la altura de vuelo: a \SI{250}{m} sobre el
sitio, la altura media de la etapa~2, sería un \SI{3}{\percent} menor y $V_2$ subiría
\SI{1.5}{\percent}; ese efecto no está incluido.

\begin{table}[H]
\centering\small
\caption{Resultados del Monte Carlo (media $\pm$ desviación; percentiles; probabilidades de cumplir).}\label{tab:arq-mc}
\begin{tabular}{@{}lccccc@{}}
\toprule
\textbf{Caso} & $V_2$ [m/s] & $V_2$ p95 & $\langle V\rangle$ p95 & $P(\text{C4, régimen})$ & $P(\text{C4, global})$\\\midrule
Arq.~1, $R=\SI{0.25}{m}$ & $6{,}60\pm0{,}42$ & 7,35 & 8,29 & 99,7\,\% & 86,0\,\%\\
Arq.~1, $R=\SI{0.23}{m}$ & $7{,}15\pm0{,}45$ & 7,94 & 8,86 & 95,9\,\% & 45,9\,\%\\
Arq.~1, $R=\SI{0.20}{m}$ & $8{,}15\pm0{,}51$ & 9,05 & 9,88 & 41,1\,\% & 0,9\,\%\\
Arq.~2, $R=\SI{0.25}{m}$ & $5{,}66\pm0{,}38$ & 6,32 & 7,28 & 100\,\% & 99,8\,\%\\
Arq.~2, $R=\SI{0.23}{m}$ & $6{,}13\pm0{,}41$ & 6,84 & 7,80 & 99,96\,\% & 97,8\,\%\\
Arq.~2, $R=\SI{0.20}{m}$ & $6{,}99\pm0{,}46$ & 7,79 & 8,72 & 97,6\,\% & 59,5\,\%\\
\bottomrule
\end{tabular}
\end{table}
\Req{C3} ($V_1$ entre 12 y 18 m/s) se cumple en el \SI{100}{\percent} de los casos.

\begin{figure}[H]
\centering
\begin{tikzpicture}
\begin{axis}[width=0.85\textwidth, height=6.2cm, grid=major, xlabel={Velocidad [m/s]}, ylabel={Densidad de probabilidad},
  xmin=3, xmax=10, ymin=0, legend style={font=\scriptsize, at={(0.98,0.98)}, anchor=north east}]
\addplot[crot, very thick] table[x=V, y=pdf_V2] {analysis/data/v2_hist_A1_R25.dat};
\addlegendentry{$V_2$, Arq.~1}
\addplot[crot, thick, dashed] table[x=V, y=pdf_Vg] {analysis/data/v2_hist_A1_R25.dat};
\addlegendentry{$\langle V\rangle$, Arq.~1}
\addplot[cfijo, very thick] table[x=V, y=pdf_V2] {analysis/data/v2_hist_A2_R25.dat};
\addlegendentry{$V_2$, Arq.~2}
\addplot[cfijo, thick, dashed] table[x=V, y=pdf_Vg] {analysis/data/v2_hist_A2_R25.dat};
\addlegendentry{$\langle V\rangle$, Arq.~2}
\addplot[cfreno, dotted, very thick] coordinates {(8,0) (8,1.2)};
\end{axis}
\end{tikzpicture}
\caption{Distribución de $V_2$ y de $\langle V\rangle$ con $R=\SI{0.25}{m}$. La línea
punteada vertical es el límite de \SI{8}{m/s}.}
\label{fig:arq-hist}
\end{figure}

La cola superior de la Arquitectura~1 en la lectura global (14\,\% de casos sobre 8~m/s) la
produce en gran parte la densidad del aire: en un sitio a \SI{800}{m} la densidad es un
\SI{9}{\percent} menor que al nivel del mar y $V_2$ sube \SI{5}{\percent}. Por eso conviene
conocer la altura del sitio (consulta~6).

%=====================================================================
\section{Decisión}\label{sec:arq-decision}

\begin{table}[H]
\centering\small
\caption{Matriz de decisión (puntaje de 1 a 5 y pesos).}\label{tab:arq-matriz}
\begin{tabular}{@{}lccc@{}}
\toprule
\textbf{Criterio} & \textbf{Peso} & \textbf{Arq.~1} & \textbf{Arq.~2}\\\midrule
\Req{C4} con margen & 0,25 & 3 & 5\\
Riesgo normativo (consulta) & 0,15 & 5 & 2\\
Complejidad mecánica & 0,10 & 3 & 4\\
Masa y energía de aterrizaje & 0,10 & 2 & 5\\
Robustez ante falla de liberación & 0,15 & 4 & 3\\
Fabricabilidad y costo & 0,10 & 3 & 4\\
Integración con antenas & 0,05 & 4 & 3\\
Validación y ensayos & 0,10 & 3 & 3\\\midrule
\textbf{Total} & 1,00 & \textbf{3,40} & \textbf{3,75}\\
\bottomrule
\end{tabular}
\end{table}

Con pesos de Dirichlet centrados en los de la tabla, la Arquitectura~2 gana en el
\SI{81}{\percent} de los casos (\SI{67}{\percent} con pesos al azar). Los puntajes son
juicios de ingeniería, no mediciones; el criterio más influyente es el riesgo normativo:
subir su peso de 0,15 a 0,24 empata las dos arquitecturas. Es decir, la decisión depende de
la respuesta de la organización más que de la técnica.

\begin{resultado}[Recomendación]
\begin{enumerate}[leftmargin=*]
  \item \textbf{Radio $R=\SI{0.25}{m}$, rotor común} (pala de \SI{206}{mm}, BV con varilla): $V_2=\SI{6.4}{m/s}$
        y $\langle V\rangle=\SI{7.1}{m/s}$ con $m_2=\SI{0.49}{kg}$.
  \item \textbf{Soltar el drogue con enclavamiento}: anillo $\to$ confirmación $\to$ pasador;
        nunca soltarlo antes de saber que las palas están libres.
  \item \textbf{Base: Arquitectura~1}, que no depende de ningún permiso. Preparar la
        Arquitectura~2 (tubo de fibra de vidrio o PVC, pestillo redundante) como mejora si la
        organización responde las consultas 3--5 en sentido favorable.
  \item \textbf{Consultar} diámetro, definición del promedio, soltado del drogue y del
        contenedor, \Req{S6} y altura del sitio, antes del PDR.
  \item \textbf{Lo que más conviene medir} es $C_R$ del rotor (ensayo E3): explica entre el 60 y
        el \SI{80}{\percent} de la varianza de $V_2$ (capítulo~\ref{cap:sens}).
\end{enumerate}
\end{resultado}
